\documentclass{beamer}
\mode<presentation>
\usepackage{amsmath,amssymb,mathtools}
\usepackage{textcomp}
\usepackage{gensymb}
\usepackage{adjustbox}
\usepackage{subcaption}
\usepackage{enumitem}
\usepackage{multicol}
\usepackage{listings}
\usepackage{url}
\usepackage{graphicx} % <-- needed for images
\def\UrlBreaks{\do\/\do-}

\usetheme{Boadilla}
\usecolortheme{lily}
\setbeamertemplate{footline}{
  \leavevmode%
  \hbox{%
  \begin{beamercolorbox}[wd=\paperwidth,ht=2ex,dp=1ex,right]{author in head/foot}%
    \insertframenumber{} / \inserttotalframenumber\hspace*{2ex}
  \end{beamercolorbox}}%
  \vskip0pt%
}
\setbeamertemplate{navigation symbols}{}

\lstset{
  frame=single,
  breaklines=true,
  columns=fullflexible,
  basicstyle=\ttfamily\tiny   % tiny font so code fits
}

\numberwithin{equation}{section}

% ---- your macros ----
\providecommand{\nCr}[2]{\,^{#1}C_{#2}}
\providecommand{\nPr}[2]{\,^{#1}P_{#2}}
\providecommand{\mbf}{\mathbf}
\providecommand{\pr}[1]{\ensuremath{\Pr\left(#1\right)}}
\providecommand{\qfunc}[1]{\ensuremath{Q\left(#1\right)}}
\providecommand{\sbrak}[1]{\ensuremath{{}\left[#1\right]}}
\providecommand{\lsbrak}[1]{\ensuremath{{}\left[#1\right.}}
\providecommand{\rsbrak}[1]{\ensuremath{\left.#1\right]}}
\providecommand{\brak}[1]{\ensuremath{\left(#1\right)}}
\providecommand{\lbrak}[1]{\ensuremath{\left(#1\right.}}
\providecommand{\rbrak}[1]{\ensuremath{\left.#1\right)}}
\providecommand{\cbrak}[1]{\ensuremath{\left\{#1\right\}}}
\providecommand{\lcbrak}[1]{\ensuremath{\left\{#1\right.}}
\providecommand{\rcbrak}[1]{\ensuremath{\left.#1\right\}}}
\theoremstyle{remark}
\newtheorem{rem}{Remark}
\newcommand{\sgn}{\mathop{\mathrm{sgn}}}
\providecommand{\abs}[1]{\left\vert#1\right\vert}
\providecommand{\res}[1]{\Res\displaylimits_{#1}}
\providecommand{\norm}[1]{\lVert#1\rVert}
\providecommand{\mtx}[1]{\mathbf{#1}}
\providecommand{\mean}[1]{E\left[ #1 \right]}
\providecommand{\fourier}{\overset{\mathcal{F}}{ \rightleftharpoons}}
\providecommand{\system}{\overset{\mathcal{H}}{ \longleftrightarrow}}
\providecommand{\dec}[2]{\ensuremath{\overset{#1}{\underset{#2}{\gtrless}}}}
\newcommand{\myvec}[1]{\ensuremath{\begin{pmatrix}#1\end{pmatrix}}}
\let\vec\mathbf

\title{Matgeo Presentation - Problem 12.82}
\author{ee25btech11063 - Vejith}

\begin{document}


\frame{\titlepage}
\begin{frame}{Question}
let $\Vec{M}$ be a 3$\times$ 3 symmetric matrix with eigenvalues 1,2 and 3. let $\Vec{N}$ be any 3$\times$ 3  matrix with real eigenvalues such that $\Vec{MN} + \Vec{NM}= 3\Vec{I}$, where $\Vec{I}$ is a 3$\times$ 3  identity matrix. Then which of the following cannot be eigenvalue\brak{s} of the matrix $\Vec{N}$ ?
\hspace{2cm}\brak{\text{ST } 2022}
\begin{enumerate}[label=\alph*)]
    \begin{multicols}{4}
  \item $\frac{1}{4}$        
  \item $\frac{3}{4}$
  \item $\frac{1}{2}$
  \item $\frac{7}{4}$  
    \end{multicols}
\end{enumerate}
\end{frame}

\begin{frame}{Solution}
    let $\lambda_1$  , $\lambda_2$ , $\lambda_3$ be the eigenvalues of $\vec{N}$\\
By Eigenvalue decomposition
\begin{align}
    \vec{M}=\vec{QDQ}^\vec{\top}
\end{align}
where 
\begin{align}
    \vec{Q}^\top\vec{Q}=\vec{Q}\vec{Q}^\top=\vec{I}\\
    \vec{D}=\text{diag}\brak{1,2 ,3}\\
    \implies \vec{D}=\vec{Q}^\top\vec{M}\vec{Q}
\end{align}
\begin{align}
    \vec{MN}+\vec{NM}=3\vec{I}\\
    \vec{Q}^\top\brak{ \vec{MN}+\vec{NM}}\vec{Q}=\vec{Q}^\top 3\vec{I}\vec{Q}\\
    \vec{Q}^\top\vec{MN}\vec{Q} + \vec{Q}^\top\vec{NM}\vec{Q}=3\vec{I}\\
  \implies  \vec{Q}^\top\vec{MN}\vec{Q}=\brak{\vec{Q}^\top\vec{M}\vec{Q}}\brak{\vec{Q}^\top\vec{N}\vec{Q}}\\
  \implies  \vec{Q}^\top\vec{NM}\vec{Q}=\brak{\vec{Q}^\top\vec{N}\vec{Q}}\brak{\vec{Q}^\top\vec{M}\vec{Q}}
\end{align}
\end{frame}

\begin{frame}{Solution}
let
\begin{align}
 \vec{Q}^\top\vec{N}\vec{Q}=\vec{N}'\\
 \implies \vec{N}=\vec{QN}'\vec{Q}^\top\\
 \implies \vec{Q}^\top\vec{MN}\vec{Q}=\vec{D}\vec{N}'\\
 \implies \vec{Q}^\top\vec{NM}\vec{Q}=\vec{N}'\vec{D}
\end{align}
From (0.11) and (0.12), (0.7) can be written as
\begin{align}
    \vec{D}\vec{N}' + \vec{N}'\vec{D}=3\vec{I}
\end{align}
let $\brak{\vec{D}\vec{N}'}_{i,j}$ be the entry of $i$th row and $j$th column of the matrix $\vec{D}\vec{N}'$.  As $\vec{D}$ is a diagonal matrix
\begin{align}
    \brak{\vec{D}\vec{N}'}_{i,j}=\vec{D}_{i,i}\vec{N}'_{i,j}
\end{align}
let $\brak{\vec{N}'\vec{D}}_{i,j}$ be the entry of $i$th row and $j$th column of the matrix $\vec{N}'\vec{D}$.  As $\vec{D}$ is a diagonal matrix
\end{frame}

\begin{frame}{Solution}
\begin{align}
    \brak{\vec{N}'\vec{D}}_{i,j}=\vec{D}_{j,j}\vec{N}'_{i,j}\\
    \implies \brak{\vec{D}_{i,i} + \vec{D}_{j,j}}\vec{N}'_{i,j}=3\brak{\vec{I}}_{i,j}
\end{align}
for $i\neq j$
\begin{align}
     \brak{\vec{D}_{i,i} + \vec{D}_{j,j}}\vec{N}'_{i,j}=0\\
   \implies  \vec{N}'_{i,j}=0
\end{align}
for $i=j$
\begin{align}
     \brak{\vec{D}_{i,i}+ \vec{D}_{j,j}}\vec{N}'_{i,j}=3\\
     \implies \vec{N}'_{i,i}=\frac{3}{2D_{i,i}}
\end{align}
from (0.3)
\begin{align}
    \implies \text{The diagonal values of } \vec{N}' \text{ are } \frac{3}{2} \text{ and } \frac{3}{4} \text{ and } \frac{1}{2}
\end{align}
\end{frame}

\begin{frame}{Solution}
In  (0.11) $\vec{Q}$ is an orthogonal matrix and $\vec{N}'$ is a diagonal matrix so we can say that (0.11) is the eigenvalue decomposition of $\vec{N}$\\
$\implies \vec{N}=\vec{QN}'\vec{Q}^\top$
\begin{align}
    \vec{N}'=\text{diag}\brak{\lambda_1,\lambda_2 ,\lambda_3}\\
    \implies \lambda_1=\frac{3}{2} \text{ and } \lambda_2=\frac{3}{4} \text{ and } \lambda_3=\frac{1}{2}
\end{align}
\end{frame}

 % --------- CODE APPENDIX ---------
\section*{Appendix: Code}

% C program
\begin{frame}[fragile]{C Code: eigen.c}
\begin{lstlisting}[language=C]
#include <stdio.h>
#include <math.h>
#include <stdlib.h>

#define N 3
#define MAX_ITERS 100
#define EPS 1e-12
void max_offdiag(double A[N][N], int *pi, int *pj) {
    double max = 0.0;
    int imax = 0, jmax = 1;
    for (int i = 0; i < N; ++i) {
        for (int j = i+1; j < N; ++j) {
            double v = fabs(A[i][j]);
            if (v > max) {
                max = v;
                imax = i; jmax = j;
            }
        }
    }
    *pi = imax; *pj = jmax;
}
void jacobi_eigensolver(double A[N][N], double eig[N]) {
    double V[N][N] = {0}; // eigenvector matrix (not used for result but kept)
    for (int i=0;i<N;i++) V[i][i]=1.0;

    double B[N][N];
    for (int i=0;i<N;i++)
        for (int j=0;j<N;j++)
            B[i][j]=A[i][j];
for (int iter = 0; iter < MAX_ITERS; ++iter) {
        int p,q;
        max_offdiag(B, &p, &q);
\end{lstlisting}
\end{frame}

\begin{frame}[fragile]{C Code: eigen.c}
\begin{lstlisting}[language=C]
        double bpq = B[p][q];
        if (fabs(bpq) < EPS) break;

        double app = B[p][p];
        double aqq = B[q][q];

        double tau = (aqq - app) / (2.0 * bpq);
        double t;
        if (tau >= 0.0) t = 1.0 / (tau + sqrt(1.0 + tau*tau));
        else t = -1.0 / (-tau + sqrt(1.0 + tau*tau));
        double c = 1.0 / sqrt(1.0 + t*t);
        double s = t * c;
        double tauPrime = s/(1.0 + c);

        // update matrix B
        double app_new = app - t * bpq;
        double aqq_new = aqq + t * bpq;
        B[p][p] = app_new;
        B[q][q] = aqq_new;
        B[p][q] = B[q][p] = 0.0;

        for (int k = 0; k < N; ++k) {
            if (k != p && k != q) {
                double bkp = B[k][p];
                double bkq = B[k][q];
                B[k][p] = B[p][k] = bkp - s*(bkq + tauPrime*bkp);
                B[k][q] = B[q][k] = bkq + s*(bkp - tauPrime*bkq);
            }
        }
        for (int k = 0; k < N; ++k) {
            double vkp = V[k][p];
            double vkq = V[k][q];
\end{lstlisting}
\end{frame}

        \begin{frame}[fragile]{C Code: eigen.c}
\begin{lstlisting}[language=C]
            V[k][p] = vkp - s*(vkq + tauPrime*vkp);
            V[k][q] = vkq + s*(vkp - tauPrime*vkq);
        }}
    for (int i = 0; i < N; ++i) eig[i] = B[i][i];
}

int main(void) {
    
    double M[N][N] = {
        {1.0, 0.0, 0.0},
        {0.0, 2.0, 0.0},
        {0.0, 0.0, 3.0}
    };
    double lambda[N];
    jacobi_eigensolver(M, lambda);

    double mu[N];
    for (int i = 0; i < N; ++i) {
        if (fabs(lambda[i]) < 1e-14) {
            fprintf(stderr, "Eigenvalue near zero; cannot compute mu = 3/(2*lambda).\n");
            return 1;}
        mu[i] = 3.0 / (2.0 * lambda[i]);}

    FILE *fp = fopen("eigen.dat", "w");
    if (!fp) {
        perror("fopen");
        return 1;}
    for (int i = 0; i < N; ++i) {
        fprintf(fp, "%.12g\n", mu[i]);
    }
    fclose(fp);
    printf("Wrote 3 eigenvalues of N to eigen.dat\n");
    return 0;}
\end{lstlisting}
\end{frame}

% Python plotting
\begin{frame}[fragile]{Python: solution.py}
\begin{lstlisting}[language=Python]
import numpy as np

# Step 1: Define symmetric matrix M (diagonal with eigenvalues 1,2,3)
M = np.diag([1.0, 2.0, 3.0])

# Step 2: Eigen decomposition of M
eigvals_M, Q = np.linalg.eigh(M)

# Step 3: Compute N' = diag(3 / (2 * λ_i))
eigvals_N = 3.0 / (2.0 * eigvals_M)
N_dash = np.diag(eigvals_N)

# Step 4: Reconstruct N = Q * N' * Q^T
N = Q @ N_dash @ Q.T

# Step 5: Verification: M*N + N*M
verify = M @ N + N @ M

# Step 6: Display results
print("Matrix M:")
print(M)

print("\nEigenvalues of M:")
print(eigvals_M)

print("\nEigenvalues of N:")
print(eigvals_N)

print("\nMatrix N:")
print(N)

print("\nVerification (M*N + N*M):")
print(verify)

\end{lstlisting}
\end{frame}
\end{document}
